\chapter{Preliminaries: Set theory and categories}
\label{chap:prelims}

\section{Naive set theory}
\label{sec:prelims:naive-set-theory}

\subsection{Equivalence relations, partitions, quotients}
\label{subsec:prelims:equivalence-relations}

\begin{definition}[Reflexive relation]
  \label{def:prelims:is-reflexive}
  \lean{Algebra0Lean.Prelims.IsReflexive}
  \leanok
  A relation $\sim$ on a set $X$ is \emph{reflexive} if $a \sim a$ for
  every $a \in X$.
\end{definition}

\begin{definition}[Symmetric relation]
  \label{def:prelims:is-symmetric}
  \lean{Algebra0Lean.Prelims.IsSymmetric}
  \leanok
  A relation $\sim$ on a set $X$ is \emph{symmetric} if $a \sim b$
  implies $b \sim a$, for all $a, b \in X$.
\end{definition}

\begin{definition}[Transitive relation]
  \label{def:prelims:is-transitive}
  \lean{Algebra0Lean.Prelims.IsTransitive}
  \leanok
  A relation $\sim$ on a set $X$ is \emph{transitive} if $a \sim b$ and
  $b \sim c$ together imply $a \sim c$, for all $a, b, c \in X$.
\end{definition}

\begin{definition}[Partition]
  \label{def:prelims:is-partition}
  \lean{Algebra0Lean.Prelims.IsPartition}
  \leanok
  A \emph{partition} of a set $X$ is a set $c$ of subsets of $X$ such
  that $\emptyset \notin c$ and every element of $X$ lies in exactly
  one member of $c$.
\end{definition}

\begin{definition}[Equivalence classes]
  \label{def:prelims:setoid-classes}
  \uses{def:prelims:is-partition}
  \lean{Algebra0Lean.Prelims.setoidClasses}
  \leanok
  Given an equivalence relation $\sim$ on $X$, the set of equivalence
  classes of $\sim$.
\end{definition}

\begin{exercise}[I.1.2]
  \label{exer:prelims:1.2}
  \uses{def:prelims:is-partition, def:prelims:setoid-classes}
  \lean{Algebra0Lean.Prelims.isPartition_setoidClasses}
  \leanok
  The equivalence classes of an equivalence relation on $X$ form a
  partition of $X$.
\end{exercise}
\begin{proof}
  \leanok
\end{proof}

\begin{exercise}[I.1.3]
  \label{exer:prelims:1.3}
  \uses{def:prelims:is-partition, def:prelims:setoid-classes}
  \lean{Algebra0Lean.Prelims.exists_setoid_of_isPartition}
  Conversely, every partition of $X$ arises as the set of equivalence
  classes of some equivalence relation on $X$.
\end{exercise}

\begin{proposition}
  \label{prop:prelims:bijective-setoid-classes}
  \uses{def:prelims:is-partition, def:prelims:setoid-classes}
  \lean{Algebra0Lean.Prelims.bijective_setoidClasses_subtype}
  The notions of equivalence relation on $X$ and partition of $X$ are
  really equivalent: \texttt{setoidClasses} is a bijection from
  equivalence relations on $X$ to partitions of $X$.
\end{proposition}

\begin{exercise}[I.1.4]
  \label{exer:prelims:1.4}
  \lean{Algebra0Lean.Prelims.card_setoid_fin_eq_bell}
  How many different equivalence relations may be defined on a set of
  $3$ elements? More generally, the number of equivalence relations on
  a set of $n$ elements is the $n$th Bell number.
\end{exercise}

\begin{exercise}[I.1.5]
  \label{exer:prelims:1.5}
  \uses{def:prelims:is-reflexive, def:prelims:is-symmetric,
    def:prelims:is-transitive}
  \lean{Algebra0Lean.Prelims.exists_reflexive_symmetric_not_transitive}
  Give an example of a relation that is reflexive and symmetric but not
  transitive.
\end{exercise}

\begin{definition}[Congruence mod $\mathbb{Z}$]
  \label{def:prelims:mod-z-relation}
  \lean{Algebra0Lean.Prelims.modZRelation}
  \leanok
  The relation on $\R$ defined by $a \sim b \iff b - a \in \Z$.
\end{definition}

\begin{definition}[Congruence mod $\mathbb{Z}$ on the plane]
  \label{def:prelims:mod-z-relation-plane}
  \uses{def:prelims:mod-z-relation}
  \lean{Algebra0Lean.Prelims.modZRelationPlane}
  \leanok
  The relation on $\R \times \R$ defined by
  $(a_1,a_2) \sim (b_1,b_2) \iff b_1-a_1 \in \Z \text{ and } b_2-a_2 \in
  \Z$.
\end{definition}

\begin{exercise}[I.1.6]
  \label{exer:prelims:1.6}
  \uses{def:prelims:mod-z-relation, def:prelims:mod-z-relation-plane}
  \lean{Algebra0Lean.Prelims.equivalence_modZRelation,
        Algebra0Lean.Prelims.equivalence_modZRelationPlane}
  The relation $a \sim b \iff b - a \in \Z$ is an equivalence relation
  on $\R$, and likewise for the corresponding relation on $\R \times
  \R$.
\end{exercise}

\section{Functions between sets}
\label{sec:prelims:functions-between-sets}

\subsection{Definition}
\label{subsec:prelims:functions:definition}

\begin{definition}[Graph of a function]
  \label{def:prelims:graph-of}
  \lean{Algebra0Lean.Prelims.graphOf}
  \leanok
  The graph of $f : X \to Y$: the subset
  $\Gamma_f = \{(a,b) \in X \times Y \mid b = f(a)\} \subseteq X \times Y$.
\end{definition}

\begin{definition}[Graph condition]
  \label{def:prelims:is-graph}
  \lean{Algebra0Lean.Prelims.IsGraph}
  \leanok
  A subset $\Gamma \subseteq X \times Y$ is (the graph of) a function
  $X \to Y$ if every $a \in X$ has exactly one $b \in Y$ with
  $(a,b) \in \Gamma$.
\end{definition}

\begin{proposition}
  \label{prop:prelims:is-graph-graph-of}
  \uses{def:prelims:graph-of, def:prelims:is-graph}
  \lean{Algebra0Lean.Prelims.isGraph_graphOf}
  The graph of a function satisfies the graph condition.
\end{proposition}

\begin{proposition}[A function really "is" its graph]
  \label{prop:prelims:bijective-graph-of}
  \uses{def:prelims:graph-of, def:prelims:is-graph,
    prop:prelims:is-graph-graph-of}
  \lean{Algebra0Lean.Prelims.bijective_graphOf}
  Taking the graph is a bijection between functions $X \to Y$ and
  subsets of $X \times Y$ satisfying the graph condition.
\end{proposition}

\subsection{Composition of functions}
\label{subsec:prelims:composition-of-functions}

\begin{proposition}
  \label{prop:prelims:comp-assoc}
  \lean{Algebra0Lean.Prelims.comp_assoc}
  Composition of functions is associative: if $f : X \to Y$,
  $g : Y \to Z$, $h : Z \to W$, then
  $h \circ (g \circ f) = (h \circ g) \circ f$.
\end{proposition}

\begin{proposition}
  \label{prop:prelims:id-comp}
  \lean{Algebra0Lean.Prelims.id_comp}
  For $f : X \to Y$, $\id_Y \circ f = f$.
\end{proposition}

\begin{proposition}
  \label{prop:prelims:comp-id}
  \lean{Algebra0Lean.Prelims.comp_id}
  For $f : X \to Y$, $f \circ \id_X = f$.
\end{proposition}

\subsection{Injections, surjections, bijections}
\label{subsec:prelims:injections-surjections-bijections}

\begin{definition}[Isomorphic]
  \label{def:prelims:isomorphic}
  \lean{Algebra0Lean.Prelims.Isomorphic}
  \leanok
  Two sets $X$, $Y$ are \emph{isomorphic} (written
  $X \overset{\sim}{=} Y$) if there is a bijection $X \to Y$.
\end{definition}

\begin{proposition}
  \label{prop:prelims:bijective-id}
  \lean{Algebra0Lean.Prelims.bijective_id}
  The identity function $\id_X : X \to X$ is a bijection.
\end{proposition}

\begin{proposition}
  \label{prop:prelims:finite-card-eq-of-isomorphic}
  \uses{def:prelims:isomorphic}
  \lean{Algebra0Lean.Prelims.finite_and_card_eq_of_isomorphic}
  If $X$ is a finite set isomorphic to $Y$, then $Y$ is finite too, and
  $|X| = |Y|$.
\end{proposition}

\begin{proposition}
  \label{prop:prelims:bijective-unit-prod-mk}
  \lean{Algebra0Lean.Prelims.bijective_unitProdMk}
  The function $f : X \to \{0\} \times X$ defined by $f(a) = (0,a)$ is
  a bijection.
\end{proposition}

\subsection{Injections, surjections, bijections: second viewpoint}
\label{subsec:prelims:injections-surjections-bijections-second-viewpoint}

\begin{proposition}[I.2.4]
  \label{prop:prelims:inv-inj-sur}
  \lean{Algebra0Lean.Prelims.hasLeftInverse_iff_injective,
        Algebra0Lean.Prelims.hasRightInverse_iff_surjective}
  \leanok
  Assume $X \neq \emptyset$, and let $f : X \to Y$ be a function.
  Then $f$ has a left-inverse if and only if it is injective, and $f$
  has a right-inverse if and only if it is surjective.
\end{proposition}
\begin{proof}
  \leanok
\end{proof}

\begin{corollary}[I.2.5]
  \label{coro:prelims:bijective-iff-inverse}
  \uses{prop:prelims:inv-inj-sur}
  \lean{Algebra0Lean.Prelims.bijective_iff_hasInverse}
  \leanok
  A function $f : X \to Y$ is a bijection if and only if it has a
  (two-sided) inverse.
\end{corollary}
\begin{proof}
  \uses{prop:prelims:inv-inj-sur}
  \leanok
\end{proof}

\begin{proposition}
  \label{prop:prelims:no-right-inverse-of-injective-not-surjective}
  \lean{Algebra0Lean.Prelims.not_hasRightInverse_of_injective_not_surjective}
  If $f$ is injective but not surjective, it has no right-inverse.
\end{proposition}

\begin{proposition}
  \label{prop:prelims:many-left-inverses}
  \lean{Algebra0Lean.Prelims.exists_ne_leftInverse_of_injective_not_surjective}
  If $f$ is injective but not surjective, and $X$ has at least two
  elements, then $f$ has more than one left-inverse.
\end{proposition}

\begin{definition}[Section]
  \label{def:prelims:is-section}
  \lean{Algebra0Lean.Prelims.IsSection}
  \leanok
  A right-inverse of $f$ is also called a \emph{section} of $f$.
\end{definition}

\begin{proposition}
  \label{prop:prelims:many-sections}
  \uses{def:prelims:is-section}
  \lean{Algebra0Lean.Prelims.exists_ne_section_of_surjective}
  If $f$ is surjective and some fiber of $f$ has at least two
  elements, then $f$ has more than one section.
\end{proposition}

\begin{definition}[Fiber]
  \label{def:prelims:fiber}
  \lean{Algebra0Lean.Prelims.fiber}
  \leanok
  The fiber of $f$ over $q \in Y$: the subset
  $f^{-1}(q) = \{a \in X \mid f(a) = q\} \subseteq X$.
\end{definition}

\begin{proposition}
  \label{prop:prelims:surjective-iff-fiber-nonempty}
  \uses{def:prelims:fiber}
  \lean{Algebra0Lean.Prelims.surjective_iff_forall_fiber_nonempty}
  $f$ is surjective if and only if every fiber of $f$ is nonempty.
\end{proposition}

\begin{proposition}
  \label{prop:prelims:injective-iff-fiber-subsingleton}
  \uses{def:prelims:fiber}
  \lean{Algebra0Lean.Prelims.injective_iff_forall_fiber_subsingleton}
  $f$ is injective if and only if every fiber of $f$ has at most one
  element.
\end{proposition}

\begin{proposition}
  \label{prop:prelims:image-inverse-eq-preimage}
  \lean{Algebra0Lean.Prelims.image_inverse_eq_preimage}
  If $g$ is a two-sided inverse of $f$, then for every $T \subseteq Y$,
  $g(T) = f^{-1}(T)$.
\end{proposition}

\subsection{Monomorphisms and epimorphisms}
\label{subsec:prelims:functions:monomorphisms-and-epimorphisms}

\begin{definition}[Monomorphism]
  \label{def:prelims:monomorphic}
  \lean{Algebra0Lean.Prelims.Monomorphic}
  \leanok
  A function $f : X \to Y$ is a \emph{monomorphism} if for every set
  $Z$ and all functions $\alpha', \alpha'' : Z \to X$,
  $f \circ \alpha' = f \circ \alpha''$ implies $\alpha' = \alpha''$.
\end{definition}

\begin{definition}[Epimorphism]
  \label{def:prelims:epimorphic}
  \lean{Algebra0Lean.Prelims.Epimorphic}
  \leanok
  A function $f : X \to Y$ is an \emph{epimorphism} if every function
  $g : Z \to Y$ factors as $g = f \circ h$ for some $h : Z \to X$.
\end{definition}

\begin{proposition}[I.2.3]
  \label{prop:prelims:2.3}
  \uses{def:prelims:monomorphic}
  \lean{Algebra0Lean.Prelims.injective_iff_monomorphic}
  \leanok
  A function is injective if and only if it is a monomorphism.
\end{proposition}
\begin{proof}
  \leanok
\end{proof}

\begin{exercise}[I.2.5]
  \label{exer:prelims:2.5}
  \uses{def:prelims:epimorphic}
  \lean{Algebra0Lean.Prelims.epimorphic_iff_surjective}
  \leanok
  A function is surjective if and only if it is an epimorphism.
\end{exercise}
\begin{proof}
  \leanok
\end{proof}

\subsection{Canonical decomposition}
\label{subsec:prelims:canonical-decomposition}

\begin{definition}[Kernel pair]
  \label{def:prelims:kernel-pair-setoid}
  \lean{Algebra0Lean.Prelims.kernelPairSetoid}
  \leanok
  The equivalence relation on $X$ induced by a function $f : X \to Y$,
  identifying $x, x'$ whenever $f(x) = f(x')$.
\end{definition}

\begin{definition}[Canonical projection]
  \label{def:prelims:canonical-projection}
  \uses{def:prelims:kernel-pair-setoid}
  \lean{Algebra0Lean.Prelims.canonicalProjection}
  \leanok
  The canonical surjection $X \twoheadrightarrow X/{\sim}$ onto the
  quotient by the kernel pair of $f$.
\end{definition}

\begin{definition}[Canonical bijection]
  \label{def:prelims:canonical-bijection}
  \uses{def:prelims:kernel-pair-setoid}
  \lean{Algebra0Lean.Prelims.canonicalBijection}
  \leanok
  The induced bijection $X/{\sim} \to \image f$.
\end{definition}

\begin{definition}[Canonical inclusion]
  \label{def:prelims:canonical-inclusion}
  \lean{Algebra0Lean.Prelims.canonicalInclusion}
  \leanok
  The inclusion $\image f \hookrightarrow Y$.
\end{definition}

\begin{theorem}[I.2.8, canonical decomposition]
  \label{thm:prelims:canonical-decomposition}
  \uses{def:prelims:canonical-projection, def:prelims:canonical-bijection,
    def:prelims:canonical-inclusion}
  \lean{Algebra0Lean.Prelims.canonicalDecomposition}
  \leanok
  Every function $f : X \to Y$ factors as
  \[
    f = (\text{inclusion}) \circ (\text{bijection}) \circ (\text{projection}),
  \]
  i.e.\ as a surjection, followed by a bijection, followed by an
  injection.
\end{theorem}
\begin{proof}
  \leanok
\end{proof}
